\newpage\thispagestyle{hl}
\hlbanner{Part I $\cdot$ Architectural foundations}{1}{The design of a decoupled scientific engine}
\hlsection{The question}
\begin{hlquote}A colleague asks for a new element and wants it in the static, modal and nonlinear solvers. Where each solver knows its elements, that means opening three solvers and adding a branch to each. Here it takes one subclass and one registry line, and no solver is touched.\end{hlquote}
\vspace{8pt}\begin{hlsum}\textbf{A new element takes one subclass and one registry line, and no solver is touched.} Material, mesh, element, system, solver and post-processing are separate layers, and a solver needs only the matrix, or the force and tangent, from an element. A twenty-line bar element runs in the existing static and modal solvers unchanged. All \cdeRegTotal{} registered elements pass an interface test, and \rom{} never imports \fea{}. The \code{auto} backend setting has a cost on sparse systems that the documentation did not mention, and the chapter measures it.\end{hlsum}
\hlsection{Key concepts}
\begin{description}\setlength\itemsep{5pt}\setlength\parsep{0pt}\renewcommand{\descriptionlabel}[1]{\hspace{\labelsep}\sffamily\bfseries\color{hlnavy}#1}
\item[Element interface] the small set of methods that every element provides to the solvers
\item[Registry] a dictionary from names to element classes, so a new element is one line
\item[Semi-discrete equation] $\bm M\ddot{\bm u}+\bm C\dot{\bm u}+\bm K\bm u=\bm F$, the one form all solvers work on
\item[Conformance test] a check that every registered element meets the interface
\item[One-way dependence] \rom{} works on plain arrays and never imports \fea{}
\item[Backend] the linear-algebra library, SciPy by default and PyTorch optionally
\end{description}
\hlsection{After this chapter you can}
\begin{itemize}\setlength\itemsep{3pt}
\item add an element by subclassing and know which interface checks it has to pass
\item say why \rom{} does not import \fea{}, and how that is checked
\item explain what the \code{auto} backend costs on sparse systems
\end{itemize}
\newpage\thispagestyle{hl}
\noindent{\sffamily\small\color{hlgrey}Chapter 1 \quad The design of a decoupled scientific engine}\par
\hlfig{\hltikz{ch01_flow}}{The flow of an analysis: material, mesh, element, system, solver. The link to \rom{} is plain arrays, never an import.}
\hlsection{Headline result}
\begin{hlheadline}\cdeRegPass{} of \cdeRegTotal{} registered elements pass the interface test, and the source of \fea{} has \cdeFeaImportsRom{} imports of \rom{} (and \rom{} has \cdeRomImportsFea{} of \fea{}).\end{hlheadline}
\hlsection{From the code}
\begin{lstlisting}[style=hlpy]
class Bar1D(Element):                         # the new element
    n_nodes, dofs_per_node, dim, gauss_order = 2, 1, 1, 1
    dof_names = ("u",)
    def shape_and_derivs(self, natural_coords):
        (xi,) = natural_coords
        return np.array([(1-xi)/2, (1+xi)/2]), np.array([[-0.5, 0.5]])
    def B_matrix(self, natural_coords, elem_coords):
        _, dN = self.shape_and_derivs(natural_coords)
        J, detJ = jacobian(dN, elem_coords)
        return np.linalg.solve(J, dN), detJ

ELEMENT_REGISTRY["bar1d"] = Bar1D             # the one registry line
\end{lstlisting}
\hlsection{A stated limit}
\begin{hllimit}The interface test is shallow: it checks attributes, shape functions and the presence of the force and tangent methods, not that an element is correct.\end{hllimit}
